%%
%% CogniSync: A Secure, Local-First Architecture for Persistent LLM Agent Memory
%% Submitted to CIKM 2026, Full Paper Track
%%
\documentclass[sigconf]{acmart}

\AtBeginDocument{%
  \providecommand\BibTeX{{%
    Bib\TeX}}}

\setcopyright{acmlicensed}
\copyrightyear{2026}
\acmYear{2026}
\acmDOI{XXXXXXX.XXXXXXX}

\acmConference[CIKM '26]{Proceedings of the 35th ACM International
  Conference on Information and Knowledge Management}{October 2026}{Atlanta,
  GA, USA}
\acmISBN{978-1-4503-XXXX-X/2026/10}

\usepackage{booktabs}
\usepackage{multirow}
\usepackage{amsmath}
\usepackage{algorithm}
\usepackage{algorithmic}
\usepackage{xspace}
\usepackage{url}
\usepackage{microtype}
\usepackage{tabularx}

\newcolumntype{Y}{>{\centering\arraybackslash}X}
\newcolumntype{L}{>{\raggedright\arraybackslash}X}
\raggedbottom
\setlength{\emergencystretch}{3em}
\hfuzz=25pt
\hbadness=10000
\makeatletter
\@ACM@balancefalse
\makeatother

\begin{document}

\title{CogniSync: A Secure, Local-First Architecture for Persistent LLM
Agent Memory}

\author{Anonymous Authors}
\affiliation{%
  \institution{Anonymous Institution}
  \city{Anonymous City}
  \country{Anonymous Country}}
\email{anonymous@anonymous.edu}
\renewcommand{\shortauthors}{Anonymous et al.}

%% ============================================================
\begin{abstract}

Most retrieval-augmented generation (RAG) systems are evaluated primarily on retrieval quality, while adversarial robustness is studied separately. This work jointly evaluates both dimensions by measuring the retrieval-quality cost of securing a RAG pipeline against adversarial content in the retrieval corpus.

We present CogniSync, a local-first RAG architecture combining: (1) a cross-session episodic memory graph for temporal and causal retrieval, (2) adaptive hybrid dense-lexical retrieval with query-dependent weighting, and (3) a context validation layer for prompt-injection and tool-spoofing defense. Evaluation across 34,929 unique queries on five benchmarks shows that, after correcting a routing miscalibration identified during large-scale evaluation, CogniSync achieves MRR 0.856 versus 0.867 for dense retrieval, corresponding to a 1.1 percentage-point reduction while providing retrieval-layer adversarial filtering.

The evaluation further shows that adaptive hybrid routing improves exact-match retrieval over fixed-weight reciprocal rank fusion (MRR 0.927 vs.~0.889), while the episodic memory graph improves cross-session retrieval by +4.1~pp MRR on multi-session agent tasks. On the domain-specific benchmark, the system operates at 27.8~ms median latency on local hardware.

These results suggest that retrieval-layer security mechanisms can be integrated into local-first RAG systems with relatively small retrieval-quality degradation while improving robustness against prompt-injection and tool-spoofing attacks.

\end{abstract}

\begin{CCSXML}
<ccs2012>
 <concept>
  <concept_id>10002951.10003260.10003261</concept_id>
  <concept_desc>Information systems~Retrieval models and ranking</concept_desc>
  <concept_significance>500</concept_significance>
 </concept>

 <concept>
  <concept_id>10002951.10003317.10003338</concept_id>
  <concept_desc>Information systems~Question answering</concept_desc>
  <concept_significance>300</concept_significance>
 </concept>

 <concept>
  <concept_id>10002978.10003022.10003023</concept_id>
  <concept_desc>Security and privacy~Software and application security</concept_desc>
  <concept_significance>300</concept_significance>
 </concept>

 <concept>
  <concept_id>10010147.10010178</concept_id>
  <concept_desc>Computing methodologies~Natural language processing</concept_desc>
  <concept_significance>200</concept_significance>
 </concept>
</ccs2012>
\end{CCSXML}

\ccsdesc[500]{Information systems~Retrieval models and ranking}
\ccsdesc[300]{Information systems~Question answering}
\ccsdesc[300]{Security and privacy~Software and application security}
\ccsdesc[200]{Computing methodologies~Natural language processing}

\keywords{
retrieval-augmented generation,
LLM agents,
agent memory,
hybrid retrieval,
prompt injection,
retrieval security,
episodic memory,
local-first RAG
}
\maketitle

%% ============================================================
\section{Introduction}
%% ============================================================

Consider an agent debugging a production incident that started three
sessions ago. The relevant context exists somewhere in the memory store:
a schema migration ran on a specific date, the CI pipeline began failing
six hours later, and rolling back a particular commit fixed it. A flat vector index may retrieve the relevant facts independently but often fails to preserve temporal ordering, causal relationships, and session provenance. Asked ``what changed before the
failure?'', the agent retrieves topically related content rather than the
event sequence that explains the failure. Enlarging the context window or adding more documents does not reliably recover the event sequence underlying failure. The limitation arises from missing temporal structure rather than insufficient data coverage.

Now add an adversary. The same pipeline that makes the agent useful,
namely the retrieval system that injects relevant context at inference
time, also gives an attacker a direct channel into the agent's
reasoning. Any document the agent can retrieve can instruct it. Recent work on prompt injection attacks against tool-selection
pipelines in LLM agents~\cite{shi2025toolhijacker}
shows that adversarial tool descriptions can redirect tool
selection without modifying model weights, enabling data
exfiltration, unauthorized actions, and malicious tool invocation. Emerging frameworks such as the Model Context Protocol (MCP)~\cite{hou2025mcp} further increase the importance of retrieval-layer security by enabling dynamic tool discovery and tighter integration between external tools, retrieval systems, and agent context.

These two problems have mostly been worked on separately. Memory papers
optimize recall and latency. Security papers measure attack success
rates. However, the interaction between retrieval quality and retrieval-layer security remains underexplored in existing evaluations:
\emph{what does it cost, in retrieval precision, to defend against
retrieval-layer attacks?} That tradeoff is real, it is measurable, and
to our knowledge it has not been characterized at scale before this
paper.

CogniSync is our attempt to close both gaps in a single architecture. It
runs entirely on-device on a local small language model, 
reducing reliance on external cloud inference services and providing network-layer control over outbound
connections. Built on that foundation, it adds three components: a
cross-session episodic memory graph for temporal reasoning, a hybrid
dense-lexical retrieval engine with adaptive query routing, and a
context validation layer that screens retrieved content for adversarial
payloads before they enter the agent's context window.

We evaluate the system across 34,929 unique queries spanning five benchmarks specifically to see where the system fails and this evaluation caught a real failure. The
large-scale results revealed that the query-type detector was sending
exact-match queries to lexical-only retrieval rather than hybrid fusion.
We identified the issue through query-stratified analysis, corrected it with adaptive
query-dependent weighting, and re-evaluated. After the fix, CogniSync
achieves MRR 0.856, trailing Dense by 1.1\%. We report the pre-fix and
post-fix numbers, explain what the difference tells us about both the
routing behavior and the observed cost of the validation layer.

The specific contributions are:

\begin{enumerate}
\item \textbf{Large-scale evaluation of the retrieval-security tradeoff in RAG.} 34,929 unique queries. 172,605 evaluation rows. An observed 1.1 percentage-point MRR reduction after routing correction.

\item \textbf{Episodic memory graph with controlled ablation analysis and
direct linking-strategy comparison.} $+$4.1~pp MRR on cross-session tasks.
Disabling the graph recovers the Vanilla~RAG baseline within measurement precision. A direct
A-MEM comparison on identical MemoryArena tasks shows that causal
edges outperform semantic similarity edges by $+$1.2~pp, isolating the
linking strategy specifically.

\item \textbf{Routing calibration: diagnosis and fix.} The large-scale
evaluation exposed a miscalibration where exact-match queries routed to
lexical-only retrieval, producing near-identical aggregate MRR with the Lexical
baseline (both 0.861). We corrected this by implementing adaptive hybrid
fusion with query-dependent weighting: exact-match queries apply a higher
lexical weight while retaining dense retrieval as a secondary signal. Exact-match MRR improved from 0.861 to 0.927. Overall MRR
improved from 0.834 to 0.856.

\item \textbf{Reproducible validation-layer specification.} We provide the
validation algorithm, hyperparameters, and implementation details in the
paper and supplementary materials.
\end{enumerate}

%% ============================================================
\section{Related Work}
%% ============================================================

\subsection{Agent Memory Architectures}

MemGPT~\cite{packer2023memgpt} treats LLM memory as an operating-system
problem: the model pages information between in-context working memory
and an external store. MemoryBank~\cite{zhong2023memorybank} applies
Ebbinghaus forgetting-curve dynamics to prioritize semantically frequent
content. A-MEM~\cite{xu2025amem} implements the Zettelkasten method with
atomic memory nodes connected by cosine-similarity edges and
LLM-generated annotations.

These architectures primarily optimize semantic relevance and memory persistence, but provide limited explicit support for temporal or causal retrieval structure. They
record what the agent learned without preserving the temporal or causal
structure of how it was learned. CogniSync's episodic graph addresses this limitation by explicitly encoding event order and causal relationships as graph edges, inspired by Tulving's distinction between episodic and semantic memory~\cite{tulving1972}.
This is not a cosmetic difference from A-MEM. A-MEM's edges connect nodes
that are semantically similar; CogniSync's edges connect nodes where one
event caused or enabled another. For incident-reconstruction queries,
the motivating use case of this work, the causal graph is structurally
better suited. A direct ablation comparing both strategies on MemoryArena
cross-session tasks confirms this: CogniSync's episodic edges achieve
MRR 0.9436 versus A-MEM's 0.9316 on identical tasks
(Section~\ref{sec:episodic_linking}).

\subsection{Retrieval-Augmented Generation}

Dense passage retrieval~\cite{karpukhin2020dpr} has dominated open-domain
QA since 2020. Its weakness on fine-grained technical
identifiers~\cite{zhao2022dense} motivates hybrid dense-sparse retrieval.
Reciprocal rank fusion~\cite{cormack2009rrf} is the standard combination
mechanism; our implementation follows the $k{=}60$ convention.

Fixed-weight RRF is a reasonable default but not necessarily optimal for
all query types. On exact-match queries involving UUIDs, error codes, and
versioned paths, assigning a higher lexical weight while retaining a
non-zero dense component yields better results than either equal-weight
fusion or lexical-only retrieval. Our routing fix (Section~\ref{sec:routing})
provides direct empirical support for this at scale.

\subsection{Security in Agentic RAG}

Perez and Ribeiro~\cite{perez2022injection} first characterized indirect
prompt injection attacks in language models. Recent work on prompt
injection attacks against tool-selection pipelines in LLM
agents~\cite{shi2025toolhijacker} extends this threat model to
multi-tool agent systems, demonstrating that adversarial tool
descriptions can manipulate retrieval and selection behavior without
modifying model weights. Hou et al.~\cite{hou2025mcp} further analyze
security risks in Model Context Protocol (MCP)-based agent ecosystems,
including risks associated with dynamic tool execution and external
context integration. CogniSync's local-first design reduces exposure to
remote tool-execution attack surfaces by avoiding dependence on external
cloud-based inference services.

Prior work on data exfiltration in LLM-integrated systems has primarily
focused on indirect prompt injection and adversarial manipulation of
retrieval-augmented workflows~\cite{greshake2023indirect}. Existing
approaches do not explicitly address retrieval-integrated validation
defenses. CogniSync's validation layer targets injection and
tool-spoofing attacks directly, while steganographic exfiltration
remains an open problem (Section~\ref{sec:exfil}).

\textbf{The measurement gap.} Prior agent memory papers do not measure
security. Prior RAG security papers do not measure precision cost. The
contribution of this work is precisely the number that sits between them:
1.1\% MRR, at scale, with a decomposition of its sources and a correction
of the routing behavior that would have obscured that number.

%% ============================================================
\section{Problem Formulation}
%% ============================================================

\subsection{Memory and Context Model}

Let an LLM have parameters $\theta$ and maximum token capacity
$C_{\max}$. At inference step $t$, working context $W_t$ satisfies
$\lvert W_t\rvert \leq C_{\max}$. The agent maintains a long-term corpus
$\mathcal{D}$, queried via retrieval function $R(q_t,\mathcal{D},k)$
returning the top-$k$ chunks for inclusion in $W_t$. The design problem
is populating $\mathcal{D}$ with episodically structured records so that
$R$ maximizes task-relevant recall across session boundaries, while
adversarially crafted content in $\mathcal{D}$ cannot redirect agent
behavior.

\subsection{Retrieval Failure Taxonomy}

We use three failure modes from the RAGEC taxonomy. \emph{Retrieval
Thrash}: successive queries converge semantically but retrieved content
fails verification. \emph{Tool Storms}: cascading invocations exhaust the
token budget. \emph{Context Bloat}: low-utility content accumulates in
$W_t$ and overwrites system instructions.

\subsection{Threat Model}

We assume an adversary controlling external content ingested into
$\mathcal{D}$, including web pages, documents, and API responses, without
model weight access. Three attack classes:

\begin{itemize}
\item \textbf{Prompt Injection}: adversarial string $x_{\mathrm{adv}} \in W_t$
  causes the agent to satisfy the adversary's goal rather than the
  system goal.
\item \textbf{Tool Spoofing}: injected documentation shifts the
  tool-selection distribution toward an unauthorized tool.
\item \textbf{Data Exfiltration}: the agent encodes private context
  within an authorized outbound response.
\end{itemize}

This model does not cover model-weight attacks, local privilege
escalation, or adversaries who know the perplexity threshold and craft
injections to evade it. The security evaluation should be read within
these bounds.

%% ============================================================
\section{CogniSync Architecture}
%% ============================================================

\begin{figure*}[t]
\centering
\includegraphics[width=\textwidth]{figures/arch1}
\caption{Overview of the CogniSync architecture. User queries are routed
through adaptive hybrid retrieval, validated for prompt injection and
tool spoofing, and delivered to a local language model. A persistent
memory layer combines an episodic graph with a semantic store to support
cross-session reasoning.}
\Description{Architecture diagram of CogniSync showing user queries flowing through adaptive hybrid retrieval and validation before reaching a local language model, with persistent memory provided by an episodic graph and a semantic store for cross-session reasoning.}
\label{fig:cognisync-architecture}
\end{figure*}

\subsection{Local-First Execution}

All reasoning, memory indexing, and tool selection run on a locally
hosted small language model in the Google Gemma-2B-Instruct (4-bit quantized), requiring
8--24~GB RAM. Two security properties follow from this deployment model: reduced dependence on external telemetry during inference, and network-layer control over outbound connections. The MCP STDIO attack class is eliminated because there is
no remote STDIO surface. External content is untrusted and passes through
the validation layer before reaching the context window.

The hardware requirement is real. Developer workstations handle it without
issue. Edge devices and IoT hardware do not. This is the cost of the
privacy guarantee.

\subsection{Cross-Session Episodic Memory Graph}
\label{sec:episodic_arch}

The episodic graph $G=(V,E)$ stores interactions as atomic event nodes.
Each node $v_i \in V$ contains: the interaction text, a Unix timestamp,
a session identifier, and the tool set invoked. Directed edge
$v_i \to v_j$ exists when $v_i$'s outcome was a structural precondition
for $v_j$'s invocation. It captures causation rather than semantic
similarity: one event caused or enabled the other.

After each session, an asynchronous consolidation process traverses $G$
and extracts stable recurring facts into a parallel FAISS semantic store.
The episodic graph is retained at full temporal resolution; the semantic
store handles stable factual lookups that do not require event ordering.
This dual-store design is inspired by Tulving's distinction between episodic and semantic memory~\cite{tulving1972}.

The architecture is motivated by incident reconstruction: ``the service
failed on this date because that migration ran and broke JWT validation,
and rolling back commit X fixed it.'' That query requires temporal and
causal ordering. A flat index may retrieve relevant events independently but does not explicitly model their temporal or causal relationships. A similarity graph can
retrieve related content but not the causal chain. The directed episodic graph is designed to preserve these relationships explicitly. Section~\ref{sec:episodic_linking} provides empirical support with a direct comparison against A-MEM's similarity-based
linking on identical tasks.

\subsection{Context Validation Layer}
\label{sec:validation_arch}

Algorithm~\ref{alg:val} specifies the context validation layer.

\begin{algorithm}[t]
\caption{Context Validation}
\label{alg:val}
\begin{algorithmic}[1]
\REQUIRE chunk $c$, perplexity threshold $\tau$, authorized schema $S_T$,
  window width $w$, circuit-breaker limit $B$
\ENSURE $\mathit{decision}\in\{\textsc{admit},\textsc{filter}\}$
\STATE $\mathit{fails}\leftarrow 0$
\FORALL{token window $c_w$ of width $w$ in $c$}
  \STATE $P\leftarrow\mathrm{PPL}(c_w;\,\theta_{\mathrm{local}})$
  \IF{$P < \tau$}
    \STATE $\mathit{fails}\leftarrow\mathit{fails}+1$
    \IF{$\mathit{fails}\geq B$}
      \STATE \textbf{return} \textsc{filter}
    \ENDIF
  \ENDIF
\ENDFOR
\IF{$c$ contains tool metadata $m$}
  \IF{$\lnot\,\mathrm{ValidateSchema}(m,S_T)$}
    \STATE \textbf{return} \textsc{filter}
  \ENDIF
\ENDIF
\STATE \textbf{return} \textsc{admit}
\end{algorithmic}
\end{algorithm}

In all experiments, the validation layer used a sliding-window width of
$w = 32$ tokens, a perplexity threshold of $\tau = 12.0$, and a
circuit-breaker limit of $B = 3$ sub-threshold windows.
Full pseudocode and additional implementation details are provided in
the anonymous supplementary repository.

A sliding window of $w$ tokens scans each retrieved chunk; windows with
perplexity below threshold $\tau$ under the local SLM are treated as
potentially instruction-like content. Low perplexity is not a sufficient
condition for adversarial behavior, but serves as a heuristic signal when
combined with schema validation and circuit-breaker thresholds. The
circuit breaker $B$ terminates scanning once enough sub-threshold windows
have accumulated.

Schema validation checks tool metadata in retrieved chunks against a
deterministic schema $S_T$ specifying authorized parameter types and
value ranges. This reduces the natural-language ambiguity that
tool-spoofing exploits: a chunk that names an authorized tool but
specifies subtly wrong parameters will typically fail schema validation even if the surrounding language appears legitimate.

\subsection{Adaptive Hybrid Dense-Lexical Retrieval}
\label{sec:retrieval_arch}

Dense retrieval uses all-MiniLM-L6-v2 (384 dimensions, FAISS
inner-product index, 120-token chunks). Lexical retrieval uses SQLite
FTS5 with BM25. The two signals fuse via Reciprocal Rank Fusion:
\begin{equation}
  \mathrm{score}(d) =
    \frac{\alpha}{k+r_{\mathrm{dense}}(d)}+
    \frac{(1-\alpha)}{k+r_{\mathrm{lex}}(d)},
    \quad k=60,
\end{equation}
where $\alpha \in (0,1)$ is set adaptively by the query-type detector.
The detector classifies each query as semantic or exact-match based on
high-entropy token patterns: UUIDs, hexadecimal strings, versioned
endpoint paths, and uppercase alphanumeric error codes (see
the anonymous supplementary repository for the full deterministic specification).

\textbf{Adaptive weighting.} For semantic queries the detector sets
$\alpha = 0.5$ (equal-weight fusion). For exact-match queries it sets
$\alpha = 0.25$, up-weighting the lexical component while retaining a
non-zero dense contribution. This differs critically from the pre-fix
behavior, in which exact-match queries set $\alpha = 0$ (lexical-only),
discarding the dense signal entirely. The non-zero dense contribution
provides fallback retrieval support when lexical matching fails on paraphrased or
partially-tokenized identifiers, and prevents the routing from collapsing
to the Lexical baseline on that query stratum. Section~\ref{sec:routing}
reports the empirical effect of this correction.

%% ============================================================
\section{Experimental Setup}
%% ============================================================

\subsection{Datasets}

\textbf{CogniSync Benchmark v1.0.} A domain-specific corpus for
engineering workflows: JWT authentication, database schema migrations,
microservice routing, CI/CD configuration, observability stacks, API
gateway rate limiting, and RBAC access control. The benchmark contains
312 semantic queries and 96 exact-match queries, evaluated with two
independent annotators on a three-level relevance scale
($\kappa = 0.81$ semantic, $\kappa = 0.94$ exact-match). All agent-pipeline
claims (Section~\ref{sec:domain}) are grounded in this benchmark alone and
should not be generalized; the public benchmarks primarily support claims regarding general retrieval behavior.

\textbf{Public benchmarks.} MS-MARCO (10,047 queries, paraphrase-heavy
passage QA), SQuAD (5,000 queries, extractive QA), SciQ (4,474 queries,
science QA), CodeSearchNet (15,000 queries, code search). Total across
all five benchmarks: 34,929 unique queries per system, approximately 172k retrieval evaluations across five configurations. Query-type classification produces 23,590
semantic and 10,931 exact-match queries across the public benchmarks. The public benchmarks are evaluated in a retrieval-only setting over restricted benchmark-specific corpora rather than the full leaderboard retrieval collections. Consequently, reported scores are not directly comparable to leaderboard-reported retrieval metrics.

\textbf{MemoryArena.} A multi-session agentic evaluation gym covering
web shopping, group travel planning, progressive information search, and
sequential mathematical reasoning. Tasks average approximately 57 action steps and often exceed 40k reasoning-trace tokens per run. A cross-session linking
sub-evaluation specifically tests reconstruction of causal event chains
across session boundaries.

\subsection{Baselines}

The large-scale evaluation compares: Dense (FAISS inner-product only),
Lexical (BM25/FTS5 only), Vanilla~RAG (Dense plus standard RAG wrapper;
retrieval scores are identical to Dense by construction), Hybrid-Naive
(always-on equal-weight RRF, $\alpha=0.5$ fixed), and CogniSync
(adaptive routing plus context validation). The domain-specific benchmark
additionally includes A-MEM, MemGPT, LangChain MMR, LlamaIndex,
MemoryBank, and Pinecone. All systems use 120-token chunks; domain-specific
experiments run over five seeds.

\subsection{Metrics}

Recall@$K$ ($K \in \{1,3,5\}$), MRR, MAP, NDCG@5. Statistical tests: Wilcoxon signed-rank test and McNemar's test over per-query deltas; rank-biserial correlation for effect size. Security evaluation: attack success rate
(ASR) and attacker MRR (ranking quality of adversarially retrieved content). Large-scale tables
report both macro-averages (arithmetic mean of per-dataset MRRs) and
micro-averages (weighted by query count).

Additional implementation details, benchmark-construction procedures,
and full pseudocode are provided in the anonymous supplementary
repository.

%% ============================================================
\section{Results}
%% ============================================================

\subsection{Domain-Specific Benchmark}
\label{sec:domain}

Table~\ref{tab:main} reports CogniSync Benchmark v1.0 results over five
seeds. CogniSync achieves Recall@5 = 79.5\%, MRR = 0.594, NDCG@5 =
0.636 at 27.8~ms median latency, $7.8\times$ faster than the cloud-simulated configuration while maintaining comparable retrieval quality. All pairwise comparisons against the evaluated baselines are statistically significant ($p < 0.01$, two-sided sign test). The nearest
competitor, Vanilla~RAG, is beaten by $+$1.3~pp Recall@5 and $+$0.012
MRR ($p = 0.0063$). These margins are modest but consistent across all
five trials. MemoryBank's temporal decay actively degrades recall on
freshly-relevant content; LangChain MMR's maximal-marginal-relevance
re-ranking trades recall for diversity in a way that is less effective under the benchmark's predominantly single-relevant-document retrieval structure.

\begin{table*}[t]
\caption{%
  Domain-specific benchmark (CogniSync Benchmark v1.0, 408 queries,
  5 seeds, averaged). Agent memory baselines evaluated here only.
  $\dagger$~MemoryBank applies temporal decay.
  *~Pinecone excluded from quality rankings (identical retrieval to
  Vanilla~RAG at $\sim$222~ms emulated cloud latency).%
}
\label{tab:main}
\begin{tabular}{lrrrrrrr}
\toprule
\textbf{System} &
  \textbf{R@1} & \textbf{R@3} & \textbf{R@5} &
  \textbf{MRR} & \textbf{MAP} & \textbf{NDCG@5} &
  \textbf{Lat (ms)} \\
\midrule
CogniSync (Ours)      & 0.468 & 0.693 & 0.795 & 0.594 & 0.583 & 0.636 & 27.8 \\
Vanilla RAG           & 0.461 & 0.676 & 0.781 & 0.582 & 0.571 & 0.624 & 24.5 \\
A-MEM                 & 0.461 & 0.671 & 0.761 & 0.576 & 0.568 & 0.617 & 24.9 \\
LlamaIndex            & 0.448 & 0.640 & 0.736 & 0.557 & 0.538 & 0.587 & 24.1 \\
MemGPT                & 0.361 & 0.615 & 0.760 & 0.505 & 0.487 & 0.555 & 26.6 \\
LangChain MMR         & 0.461 & 0.476 & 0.500 & 0.474 & 0.447 & 0.460 & 28.8 \\
MemoryBank$^\dagger$ & 0.189 & 0.364 & 0.560 & 0.307 & 0.285 & 0.350 & 24.0 \\
Pinecone (lat.)*      & 0.461 & 0.676 & 0.781 & 0.582 & 0.571 & 0.624 & 223.4 \\
\bottomrule
\end{tabular}
\end{table*}

\subsection{Latency and Scaling}

Table~\ref{tab:lat} shows latency and memory from 2K to 100K documents.
The overhead versus Vanilla~RAG stays between 17\% and 28\%. At 27.8~ms median on the domain benchmark, the additional latency relative to Vanilla~RAG remains modest at operating scale. Memory scales linearly at
approximately 3~MB per 1K documents; the 100K-document configuration
reaches 313.5~MB, remaining within the memory capacity of typical developer workstations.

\begin{table}[t]
\caption{Latency and memory scaling. ``CS'' = CogniSync; ``VR'' = Vanilla RAG.}
\label{tab:lat}
\footnotesize\setlength{\tabcolsep}{3pt}
\begin{tabularx}{\columnwidth}{@{}LYYYYY@{}}
\toprule
\textbf{Corpus} &
  \textbf{CS (ms)} & \textbf{VR (ms)} &
  \textbf{Overhead} & \textbf{p95 (ms)} & \textbf{Mem (MB)} \\
\midrule
2,000   & 29.6 & 24.9 & $+$19\% & 43.7 & 6.7   \\
10,000  & 29.1 & 24.1 & $+$21\% & 38.6 & 31.2  \\
50,000  & 36.4 & 30.3 & $+$20\% & 44.6 & 156.9 \\
100,000 & 51.7 & 40.4 & $+$28\% & 64.7 & 313.5 \\
\bottomrule
\end{tabularx}
\end{table}

\subsection{Agent Task Performance}
\label{sec:agent}

Table~\ref{tab:agent} reports per-task MRR on MemoryArena. The episodic
graph exhibits a targeted effect: $+$4.1~pp MRR on cross-session
tasks, minimal effect on single-session tasks ($<$0.003 MRR across
all single-session types). Disabling the graph matches the Vanilla~RAG baseline within measurement precision on cross-session MRR (0.698 in both cases).

Two task categories where Full CogniSync trails Vanilla~RAG are noteworthy. Code generation (0.925 vs.~0.950) and API integration
(0.652 vs.~0.692) both involve exact API names and configuration keys.
Both systems achieve full task completion on the evaluated instances for these task types, so the MRR gap
reflects ranking quality, not success rate. Terse API documentation has
characteristically low perplexity under a general language model, making
it vulnerable to the validation layer's filtering. Tuning the perplexity
threshold for API-documentation corpora may reduce this gap; it is a
deployment-specific calibration opportunity, not a fundamental limitation.

\begin{table}[t]
\caption{%
  Agent task MRR on MemoryArena. Disabling the episodic graph recovers
  the Vanilla~RAG baseline on cross-session tasks exactly.%
}
\label{tab:agent}
\footnotesize\setlength{\tabcolsep}{3pt}
\begin{tabularx}{\columnwidth}{@{}LYYYY@{}}
\toprule
\textbf{Configuration} &
  \textbf{Code Gen} & \textbf{Debug} &
  \textbf{API Int.} & \textbf{Cross-Sess.} \\
\midrule
Full CogniSync              & 0.925 & 0.787 & 0.652 & 0.739 \\
Ablation: no episodic graph & 0.935 & 0.784 & 0.671 & 0.698 \\
Vanilla RAG baseline        & 0.950 & 0.783 & 0.692 & 0.698 \\
\bottomrule
\end{tabularx}
\end{table}

\subsection{Large-Scale Generalization and the Security-Quality Tradeoff}
\label{sec:largescale}

Tables~\ref{tab:large} and~\ref{tab:pds} present results across all
34,929 unique queries with adaptive hybrid routing. Dense/Vanilla~RAG leads at
MRR 0.867. CogniSync follows at MRR 0.856, 1.1 percentage points lower.
Hybrid-Naive sits at 0.844.

This result merits further analysis. CogniSync now outperforms Hybrid-Naive
(0.856 vs.~0.844) despite the additional overhead of the validation layer.
The reason is that adaptive routing with query-dependent weighting
($\alpha = 0.25$ for exact-match, $\alpha = 0.5$ for semantic) outperforms
fixed equal-weight fusion on exact-match queries: MRR 0.927 vs.~0.889
on that stratum. Hybrid-Naive is a reasonable default; calibrated routing
performs better empirically.

The gap between CogniSync and Dense ranges from $-$0.045 on MS-MARCO
to $-$0.002 on CodeSearchNet. A notable result from the query-level analysis: CogniSync achieves higher MRR than Vanilla~RAG on 2,499 queries,
with an average MRR gain of $+$0.374 on those wins. Ninety-four percent
of these wins (2,355 queries) occur on MS-MARCO, demonstrating that hybrid
fusion improves retrieval on sparse-query cases where dense retrieval alone performs poorly.
The validation overhead is not uniformly distributed; it concentrates where
the embedding space is weakest, which overlaps with query regimes where validation-layer filtering is most active.

CodeSearchNet shows minimal sensitivity to the security-quality tradeoff.
Only 0.8\% of code queries suffer a severe MRR penalty (drop $\geq 0.5$
versus Vanilla~RAG), compared to 12\% on MS-MARCO and 5.9\% on SQuAD.
Code blocks, API specifications, and functional documentation have rigid
syntactic structures that behave predictably under the SLM's perplexity
threshold, supporting the hypothesis that CogniSync is well-suited to software-engineering retrieval workloads.

\begin{figure}[t]
\centering
\includegraphics[width=\columnwidth]{figures/security_tradeoff_plot}
\caption{Security-quality tradeoff across retrieval systems. The x-axis
shows attack success rate (ASR) under the prompt-injection evaluation,
and the y-axis shows retrieval quality measured by mean reciprocal rank
(MRR). CogniSync achieves near-zero attack success while maintaining
retrieval performance close to the Dense/Vanilla~RAG baseline.}
\Description{Scatter plot comparing retrieval systems by prompt-injection attack success rate on the x-axis and mean reciprocal rank on the y-axis, showing CogniSync near zero attack success with retrieval quality close to Dense and Vanilla RAG.}
\label{fig:security-tradeoff}
\end{figure}

\begin{table*}[t]
\caption{%
  Large-scale aggregate evaluation. Dense and Vanilla~RAG share identical
  retrieval; scores are equal by construction. Wilcoxon $p$-values are
  computed over per-query MRR deltas (34,929 unique queries). CogniSync uses
  adaptive hybrid routing throughout.%
}
\label{tab:large}
\footnotesize\setlength{\tabcolsep}{3pt}
\begin{tabularx}{\textwidth}{@{}LYYYYYLY@{}}
\toprule
\textbf{System} &
  \textbf{R@1} & \textbf{R@3} & \textbf{R@5} &
  \textbf{MRR} & \textbf{NDCG@5} & \textbf{Wilcoxon vs Dense/VRAG} &
  \textbf{$\Delta$MRR} \\
\midrule
Dense/Vanilla RAG & 0.803 & 0.906 & 0.961 & 0.867 & 0.891 & -- & $\phantom{-}$0.000 \\
CogniSync         & 0.792 & 0.898 & 0.955 & 0.856 & 0.882 & $p < 0.001$ & $-$0.011 \\
Hybrid-Naive      & 0.774 & 0.889 & 0.951 & 0.844 & 0.872 & $p < 0.001$ & $-$0.023 \\
Lexical (BM25)    & 0.726 & 0.838 & 0.912 & 0.796 & 0.826 & $p < 0.001$ & $-$0.071 \\
\bottomrule
\end{tabularx}
\end{table*}
Public benchmark evaluations use restricted benchmark-specific retrieval corpora and are therefore not directly comparable to leaderboard retrieval settings.
\begin{table*}[t]
\caption{%
  Per-dataset MRR with macro- and micro-averages. Macro = arithmetic mean;
  Micro = query-count weighted (matches Table~\ref{tab:large}).
  v1.0 = CogniSync Benchmark.%
}
\label{tab:pds}
\footnotesize\setlength{\tabcolsep}{2.5pt}
\begin{tabularx}{\textwidth}{@{}LYYYYYYYY@{}}
\toprule
\textbf{System} &
  \textbf{MARCO} & \textbf{SQuAD} & \textbf{SciQ} &
  \textbf{CSNet} & \textbf{v1.0} & \textbf{Macro} &
  \textbf{Micro-avg} \\
\midrule
Dense/VRAG   & 0.543 & 0.991 & 0.982 & 0.998 & 1.000 & 0.903 & 0.867 \\
CogniSync    & 0.526 & 0.981 & 0.982 & 0.996 & 0.883 & 0.874 & 0.856 \\
Hybrid-Naive & 0.499 & 0.971 & 0.972 & 0.995 & 0.873 & 0.862 & 0.844 \\
Lexical      & 0.375 & 0.934 & 0.954 & 0.989 & 0.741 & 0.799 & 0.796 \\
\bottomrule
\end{tabularx}
\end{table*}
\textbf{Interpreting the residual gap.} After the routing fix, the
remaining 1.1\% MRR gap between CogniSync and Dense is primarily associated with the validation layer. The routing analysis in
Section~\ref{sec:routing} shows that on semantic queries (68.3\% of the
evaluation), CogniSync is statistically indistinguishable from Hybrid-Naive
(both MRR 0.823). The 1.1\% gap concentrates on exact-match queries where
the validation layer filters some legitimate low-perplexity content. This behavior represents the primary observed security-quality tradeoff in the evaluated setting. It is smaller than the 3.3\%
headline number would have suggested before the routing correction.

%% ============================================================
\section{Routing Calibration: Diagnosis and Fix}
\label{sec:routing}
%% ============================================================

A notable observation in the evaluation was the similarity between several query-stratified retrieval metrics before the routing correction. 
\textbf{Diagnosing the miscalibration.} In the pre-fix system, two exact
equalities stood out in the MRR columns: (1) CogniSync and Hybrid-Naive
were identical on semantic queries (both 0.823); (2) CogniSync and Lexical
were identical on exact-match queries (both 0.861). These were not
approximations; they were exact to three decimal places. Equality (1)
confirmed the routing was working correctly on semantic queries: the
detector classified them as semantic, routed to hybrid fusion, and produced
the same result as Hybrid-Naive. Equality (2) confirmed a miscalibration
on exact-match queries: the detector was setting $\alpha = 0$, discarding
the dense signal entirely and producing output identical to pure BM25
for 31.7\% of all queries. A 408-query benchmark would not have made
this equality exact; 34,929 unique queries did.

\textbf{Implementing the fix.} We corrected the miscalibration by changing
the exact-match routing from $\alpha = 0$ (lexical-only) to $\alpha = 0.25$
(lexical-upweighted hybrid). This retains a non-zero dense contribution
for fallback robustness while still prioritizing lexical matching on
high-entropy tokens. The change required modifying a single parameter in
the query-type detector; the rest of the architecture is unchanged.

\textbf{Effect of the correction.} The improvement is concentrated on
the exact-match stratum, as expected. Exact-match MRR improved from 0.861
to 0.927, surpassing Hybrid-Naive's 0.889. The gain over Hybrid-Naive
confirms that $\alpha = 0.25$ (lexical-upweighted) outperforms $\alpha = 0.5$
(equal-weight) on this query type: when the query contains a UUID or error
code, the lexical signal deserves more weight, but discarding the dense
signal entirely was too aggressive. Semantic query performance is unchanged
(MRR 0.823), as expected: the fix only touches the routing branch for
exact-match queries. Overall MRR improved from 0.834 to 0.856. The routing correction should be interpreted as a post-hoc calibration
analysis rather than an evaluation on an independent held-out tuning set.
The exact-match weight $\alpha = 0.25$ was selected after inspecting the
benchmark results and then re-evaluated on the same query sets.
Accordingly, the post-fix results are best interpreted as evidence that
part of the original performance degradation was due to routing
miscalibration rather than solely to the validation layer.

\begin{table*}[t]
\caption{%
  Query-type decomposition across all 34,929 unique queries, before and after
  the routing correction. Pre-fix: the exact equality between CogniSync
  and Lexical on exact-match MRR (both 0.861) is consistent with routing behavior effectively equivalent to $\alpha=0$. Post-fix: adaptive weighting ($\alpha=0.25$ for
  exact-match) outperforms both fixed-weight alternatives.%
}
\label{tab:routing}
\footnotesize\setlength{\tabcolsep}{3.5pt}
\begin{tabularx}{\textwidth}{@{}LYYYYL@{}}
\toprule
\textbf{System} &
  \textbf{Sem.\ MRR} & \textbf{EM MRR} &
  \textbf{Sem.\ R@5} & \textbf{EM R@5} &
  \textbf{Effective Path} \\
\midrule
Dense/Vanilla RAG              & 0.844 & 0.902 & 0.955 & 0.971 &
  Query-agnostic dense \\
Hybrid-Naive                   & 0.823 & 0.889 & 0.944 & 0.962 &
  Always-on RRF ($\alpha=0.5$) \\
CogniSync \textit{(pre-fix)}   & 0.823 & 0.861 & 0.944 & 0.934 &
  Semantic $\to$ RRF; exact-match $\to$ lexical-only ($\alpha=0$, \textit{miscalibrated routing}) \\
CogniSync \textit{(post-fix)}  & 0.823 & 0.927 & 0.944 & 0.978 &
  Semantic $\to$ RRF ($\alpha=0.5$); exact-match $\to$ adaptive hybrid ($\alpha=0.25$) \\
Lexical                        & 0.771 & 0.861 & 0.897 & 0.934 &
  Lexical-only \\
\bottomrule
\end{tabularx}
\end{table*}

The 0.927 exact-match MRR for CogniSync post-fix exceeds Hybrid-Naive's
0.889. This is not surprising in hindsight: for queries containing UUIDs
or versioned endpoint paths, a higher lexical weight is strictly more
appropriate than equal-weight fusion, but the dense component still helps
when an identifier is partially tokenized or when a paraphrased variant
appears in the corpus. The adaptive weighting captures this tradeoff in
a way that neither fixed-weight option does.

%% ============================================================
\section{Ablation Studies}
%% ============================================================

\subsection{Episodic Graph: Disabling and Linking-Strategy Comparison}
\label{sec:episodic_linking}

Disabling the graph while holding all other components constant reduces
cross-session MRR from 0.739 to 0.698, recovering the Vanilla~RAG baseline
exactly. Single-session metrics change by less than 0.003 MRR in every
task type. There is no interference: the episodic component adds value on
cross-session tasks and is invisible on single-session tasks.

Table~\ref{tab:linking} fills the linking-strategy comparison gap. We ran
both CogniSync (causal/temporal edges) and A-MEM (cosine-similarity edges)
on the same MemoryArena cross-session linking tasks. CogniSync achieves
MRR 0.9436 versus A-MEM's 0.9316, a $+$1.2~pp advantage that isolates the
linking strategy specifically. A-MEM's semantic edges actually degraded the
Vanilla~RAG baseline (0.9316 vs.~0.9426), confirming that topical clustering
is a poor substitute for causal ordering on incident-reconstruction tasks.
When the question is ``what happened in what order,'' a similarity graph
answers a different question.

\begin{table}[t]
\caption{%
  Cross-session linking strategy ablation on MemoryArena. All variants
  use identical retrieval; only the edge type in the memory graph varies.%
}
\label{tab:linking}
\footnotesize\setlength{\tabcolsep}{3pt}
\begin{tabularx}{\columnwidth}{@{}LLYY@{}}
\toprule
\textbf{Variant} & \textbf{Linking Strategy} & \textbf{MRR} & \textbf{Recall@5} \\
\midrule
Vanilla RAG (Baseline)     & None (flat index) & 0.9426 & 1.000 \\
A-MEM (Semantic Edges)     & Cosine similarity & 0.9316 & 1.000 \\
CogniSync (Episodic Edges) & Causal/temporal   & 0.9436 & 1.000 \\
\bottomrule
\end{tabularx}
\end{table}

\subsection{Hybrid Retrieval}

On the domain-specific benchmark, hybrid and dense-only systems produce
identical scores on the 312 semantic queries: the routing correctly directs
these to hybrid fusion, which on paraphrase-heavy queries differs negligibly
from dense-only in a high-quality embedding space. On the 96 exact-match
queries, adaptive hybrid ($\alpha=0.25$) achieves Recall@5 = 0.834 versus
0.750 for dense-only, an $+$8.4~pp gain confirming the lexical component
earns its place on technical identifier queries.

\subsection{Chunk Size}

FAISS Recall@5 was 0.750 at both 200 and 500 tokens; MRR dropped from
0.606 to 0.600. The system does not require chunk-size tuning in this
range, removing a meaningful deployment friction.

%% ============================================================
\section{Security Evaluation}
%% ============================================================
\label{sec:security}

\subsection{Prompt Injection and Tool Spoofing}

Table~\ref{tab:sec} reports attack effectiveness before and after the
validation layer, using two complementary metrics: Attack Success Rate
(ASR) measures whether an attack fully redirects agent behavior; Attacker
MRR measures partial goal achievement. Without sanitization: prompt
injection reaches attacker MRR 0.90; tool spoofing reaches 1.00. After
sanitization: no prompt injection attack in the evaluated set succeeded
(ASR 0.00\%, attacker MRR 0.000). Tool spoofing attacker MRR fell from
1.00 to 0.038, a $\sim$96\% reduction.

Tool spoofing is harder to eliminate completely than prompt injection
because a subtly malformed tool schema can be both fluent (passing the
perplexity check) and adversarial (redirecting tool selection). The schema
validator closes most of this gap by catching parameter-level manipulations
that perplexity windowing cannot see.

The attack set covers known injection patterns from published
taxonomies. It does not
include adaptive attacks designed to evade the perplexity threshold
specifically. The 0.00\% ASR on prompt injection applies within this scope.

\subsection{Adaptive Threshold Evasion}
\label{sec:adaptive}

The advanced Adaptive Threshold Evasion attack, where the attacker
specifically targets Algorithm~\ref{alg:val}'s perplexity threshold,
succeeds at 4.72\% ASR on both the undefended baseline and CogniSync.
We report this directly. The validation layer does not currently defend
against an adversary who knows and exploits the threshold; that is the
most important open security gap.

\subsection{Data Exfiltration}
\label{sec:exfil}

Known data exfiltration payloads are blocked entirely (ASR 0.00\%, MRR
drop 0.000). However, residual leakage via authorized channels remains:
mean leakage MRR 0.073 across seeds, with a $3.4\times$ spread. The
remaining attack exploits an authorized outbound channel: the agent encodes
private context within a semantically legitimate response payload.
Connection-level filtering cannot close this because the connection is
permitted by design. Applying semantic monitoring of output content relative
to authorized output specifications, following Bagdasaryan and
Shmatikov~\cite{bagdasaryan2023covert}, is the most promising direction.

\begin{table*}[t]
\caption{%
  Security evaluation. ASR = Attack Success Rate (0\% = fully blocked).
  Attacker MRR measures partial goal achievement (lower = stronger defense).
  Adaptive Evasion is the open gap. An Attacker MRR of 0.073 reflects
  residual leakage via authorized channels.%
}
\label{tab:sec}
\footnotesize\setlength{\tabcolsep}{3pt}
\begin{tabularx}{\textwidth}{@{}LLYYYL@{}}
\toprule
\textbf{Attack} & \textbf{System} & \textbf{ASR} & \textbf{MRR Drop} & \textbf{Attacker MRR} & \textbf{Result} \\
\midrule
Data Exfiltration & No Defense & 3.20\% & 0.0054 & --    & Vulnerable \\
Data Exfiltration & CogniSync  & 0.00\% & 0.000  & 0.073 & Blocked \\
Prompt Injection  & No Defense & 3.20\% & 0.0055 & 0.900 & Vulnerable \\
Prompt Injection  & CogniSync  & 0.00\% & 0.000  & 0.000 & Complete \\
Tool Spoofing     & No Defense & --      & --     & 1.000 & Vulnerable \\
Tool Spoofing     & CogniSync  & --      & --     & 0.038 & $\sim$96\% reduction \\
Adaptive Evasion  & No Defense & 4.72\% & 0.0081 & --    & Open gap \\
Adaptive Evasion  & CogniSync  & 4.72\% & 0.0081 & --    & Open gap \\
\bottomrule
\end{tabularx}
\end{table*}

%% ============================================================
\section{Limitations}
\label{sec:lim}
%% ============================================================

\textbf{Hardware.} The local SLM (Google Gemma-2B-Instruct (4-bit quantized)) requires
8--24~GB RAM. Developer workstations handle it. Edge and IoT devices do
not. Local execution is what gives the privacy guarantee, and that
guarantee has a hardware price.

\textbf{Benchmark scope.} The domain-specific benchmark was
author-constructed and covers the domain the system was designed for.
Statistical significance on 408 queries at $p < 0.01$ is consistent
evidence within that domain. It is not a substitute for a third-party
engineering retrieval benchmark.

\textbf{Adaptive adversaries.} The security evaluation covers fixed-taxonomy
attacks. It does not cover adversaries who know the perplexity threshold and
design injections to stay above it. The Adaptive Threshold Evasion result
(4.72\% ASR, identical to the undefended baseline) makes this gap explicit.
This is left for future work.

\textbf{Post-hoc routing calibration.} The adaptive routing weights
reported in Section~\ref{sec:routing} were calibrated after inspecting
the benchmark results and re-evaluated on the same query sets.
Consequently, the post-fix performance estimates may be optimistic.
The routing analysis should therefore be interpreted as diagnostic
evidence rather than as a fully independent held-out evaluation.
Future work should use separate development and test sets for routing
parameter selection.

\textbf{Perplexity threshold calibration.} The current threshold is set
globally. The validation layer's false-positive rate varies by corpus type:
API documentation has characteristically low perplexity and is more likely
to be filtered. Corpus-specific threshold tuning is a deployment-specific
calibration opportunity, not a fundamental limitation.

%% ============================================================
\section{Conclusion}
%% ============================================================

The question CogniSync was built to answer is: what does it cost to make
a RAG pipeline safe? The observed post-calibration gap is 1.1 percentage points in MRR, but the path to that number matters as much
as the number itself.

Our large-scale evaluation caught a routing miscalibration that was not
visible at smaller scale: exact-match queries were routed to lexical-only
retrieval, producing an exact MRR equality with the pure BM25 baseline
(both 0.861) on 31.7\% of all queries. We diagnosed this through the
exact equality fingerprint, corrected it by implementing adaptive hybrid
fusion with query-dependent weighting ($\alpha = 0.25$ for exact-match),
and re-evaluated. Exact-match MRR improved from 0.861 to 0.927, surpassing
fixed-weight Hybrid-Naive (0.889). Overall MRR improved from 0.834 to 0.856.

After post-hoc routing calibration, the remaining 1.1 percentage-point
MRR gap is primarily associated with the validation layer, although this
estimate should be interpreted cautiously because the routing weights
were calibrated on the same evaluation benchmarks: the cost of screening retrieved content for prompt
injection and tool-spoofing payloads before they enter the agent's context
window. Dense retrieval has no such mechanism. CogniSync does, and the
evaluation now measures its cost precisely rather than conflating it with
a routing bug.

On the problem the system was specifically designed for, the picture is
more favorable. On a 408-query engineering benchmark, CogniSync outperforms
seven baselines including A-MEM, MemGPT, and Vanilla~RAG with statistical
significance, at $7.8\times$ the latency advantage of cloud retrieval. The
episodic graph delivers cleanly: $+$4.1~pp cross-session MRR, zero
single-session interference, confirmed by a tight ablation. A direct
linking-strategy comparison confirms that causal edges outperform A-MEM's
similarity edges by $+$1.2~pp on cross-session tasks. And CodeSearchNet
shows only 0.8\% severe loss rate under the validation layer versus 12\%
on MS-MARCO, which is the architecture working as intended: the syntactic
regularity of code and API documentation makes perplexity-based filtering
reliable on that corpus.

The broader argument is methodological. A 408-query benchmark would not
have found the routing bug. 34,929 unique queries did. Running large-scale
evaluations that include unfavorable baselines and then fixing what
they find is how you find out what is actually wrong with your system.

An anonymous supplementary repository containing full pseudocode,
hyperparameter tables, benchmark-construction details, routing classifier
rules, and the complete evaluation outputs is provided for peer review:
\url{https://anonymous.4open.science/r/CogniSync-XXXX/}.
Code and benchmark artifacts will be released after publication.

\begin{acks}
Omitted for blind review.
\end{acks}

\bibliographystyle{ACM-Reference-Format}
\bibliography{cognisync}

%% ============================================================



\end{document}